\documentclass[11pt]{article}
\usepackage[margin=1in]{geometry}
\usepackage{xcolor}
\usepackage{tcolorbox}
\usepackage{hyperref}
\usepackage{enumitem}
\usepackage{booktabs}
\usepackage{graphicx}

% Define colors
\definecolor{osaorange}{RGB}{220,115,38}
\definecolor{aiblue}{RGB}{30,144,255}

% Configure hyperref
\hypersetup{
    colorlinks=true,
    linkcolor=blue,
    urlcolor=blue,
    citecolor=blue
}

\title{\textbf{H524 Introduction to Biostatistics\\
\Large Public Health Data Analysis Project\\
Fall 2025}}
\author{Dr. John Molitor}
\date{}

\begin{document}
\maketitle

\begin{tcolorbox}[colback=aiblue!10,colframe=aiblue,boxrule=2pt,arc=2pt,left=6pt,right=6pt,top=6pt,bottom=6pt]
\centering
\textbf{\large Final Assessment Project Overview}\\[0.3em]
Teams of 7 students will conduct a comprehensive biostatistical analysis of a real public health dataset as the culminating course assessment (25\% of final grade). This project replaces traditional final exams and demonstrates biostatistical analysis skills while maintaining scientific rigor and reproducibility.
\end{tcolorbox}

\section{Project Timeline}
\begin{itemize}[leftmargin=20pt]
\item \textbf{Week 3:} Project introduction and team formation
\item \textbf{Week 4:} Dataset selection and research question formulation
\item \textbf{Week 5:} Project proposal due (1 page)
\item \textbf{Week 7:} Mid-project check-in and feedback
\item \textbf{Week 9:} Draft analysis due for peer review
\item \textbf{Week 10:} Final project submission
\item \textbf{Finals Week:} Final assessment presentations (20 minutes + 5 minutes Q\&A per team)
\end{itemize}

\section{Available Datasets}

Each team must select one of the following real public health datasets:

\subsection{Option 1: CDC PLACES Health Data (County-Level)}
\begin{itemize}
\item \textbf{Source:} CDC PLACES Project 2022
\item \textbf{Focus Areas:} Chronic disease prevalence, health behaviors, prevention measures
\item \textbf{Sample Size:} ~3,100 counties/county equivalents
\item \textbf{Access:} \url{https://chronicdata.cdc.gov/browse?category=500+Cities+%26+Places}
\item \textbf{Data Type:} Aggregated county-level estimates (no survey weights needed)
\item \textbf{Suggested Research Questions:}
    \begin{itemize}
    \item Geographic patterns in diabetes prevalence across US counties
    \item Association between preventive care access and health outcomes
    \item Rural vs urban differences in chronic disease burden
    \end{itemize}
\end{itemize}

\subsection{Option 2: COVID-19 Case Surveillance Public Use Data}
\begin{itemize}
\item \textbf{Source:} CDC COVID-19 Case Surveillance
\item \textbf{Focus Areas:} COVID-19 outcomes, risk factors, demographics
\item \textbf{Sample Size:} ~1 million de-identified cases
\item \textbf{Access:} \url{https://data.cdc.gov/Case-Surveillance/}
\item \textbf{Suggested Research Questions:}
    \begin{itemize}
    \item Risk factors for severe COVID-19 outcomes
    \item Geographic and temporal patterns in case fatality rates
    \item Age-stratified analysis of hospitalization predictors
    \end{itemize}
\end{itemize}

\subsection{Option 3: Heart Disease Mortality Data (County-Level)}
\begin{itemize}
\item \textbf{Source:} CDC WONDER Underlying Cause of Death
\item \textbf{Focus Areas:} Heart disease mortality rates, demographic patterns
\item \textbf{Sample Size:} ~3,100 counties with 5+ years of data
\item \textbf{Access:} \url{https://wonder.cdc.gov/controller/datarequest/D77}
\item \textbf{Data Type:} Aggregated mortality rates (no survey weights needed)
\item \textbf{Suggested Research Questions:}
    \begin{itemize}
    \item Geographic disparities in heart disease mortality
    \item Age and gender patterns in cardiovascular death rates
    \item Trends in heart disease mortality over time
    \end{itemize}
\end{itemize}

\subsection{Option 4: Global Burden of Disease (GBD) Study}
\begin{itemize}
\item \textbf{Source:} IHME GBD 2019
\item \textbf{Focus Areas:} Disease burden, mortality, risk factors globally
\item \textbf{Sample Size:} 204 countries and territories
\item \textbf{Access:} \url{http://ghdx.healthdata.org/gbd-results-tool}
\item \textbf{Suggested Research Questions:}
    \begin{itemize}
    \item Trends in disease burden across income levels
    \item Environmental risk factors and respiratory disease mortality
    \item Progress toward health-related SDG targets
    \end{itemize}
\end{itemize}

\subsection{Option 5: Youth Tobacco Use Data}
\begin{itemize}
\item \textbf{Source:} CDC Youth Risk Behavior Surveillance System (State-Level)
\item \textbf{Focus Areas:} Tobacco use patterns, demographics, risk behaviors
\item \textbf{Sample Size:} 50 states + DC, multiple years
\item \textbf{Access:} \url{https://www.cdc.gov/healthyyouth/data/yrbs/}
\item \textbf{Data Type:} State-level aggregated percentages (no survey weights needed)
\item \textbf{Suggested Research Questions:}
    \begin{itemize}
    \item State-level variations in youth tobacco use
    \item Trends in e-cigarette vs traditional cigarette use
    \item Association between tobacco policies and use rates
    \end{itemize}
\end{itemize}

\subsection{Option 6: OpenFDA Adverse Events Database}
\begin{itemize}
\item \textbf{Source:} FDA Adverse Event Reporting System (FAERS)
\item \textbf{Focus Areas:} Drug safety, adverse events, pharmacovigilance
\item \textbf{Sample Size:} Millions of reports
\item \textbf{Access:} \url{https://open.fda.gov/data/faers/}
\item \textbf{Suggested Research Questions:}
    \begin{itemize}
    \item Patterns in adverse events for specific drug classes
    \item Age and gender differences in drug safety profiles
    \item Time trends in reporting of serious adverse events
    \end{itemize}
\end{itemize}

\section{Project Components}

\subsection{Part 1: Research Question and Literature Review (15\%)}
\begin{itemize}
\item Formulate a clear, answerable research question
\item Identify gaps your analysis will address
\end{itemize}

\subsection{Part 2: Data Management and Exploration (20\%)}
\begin{itemize}
\item Exploratory data analysis with appropriate visualizations
\item Missing data assessment and handling strategy
\end{itemize}

\subsection{Part 3: Statistical Analysis (30\%)}
\begin{itemize}
\item Apply appropriate statistical methods from H524 curriculum:
    \begin{itemize}
    \item Descriptive statistics and data visualization
    \item Hypothesis testing (t-tests, chi-square, ANOVA)
    \item Correlation and regression analysis
    \item Appropriate adjustments for multiple comparisons
    \end{itemize}
\item Conduct sensitivity analyses where appropriate
\end{itemize}

\subsection{Part 5: Final Presentation and Communication (25\%)}
\begin{itemize}
\item 20-minute team presentation + 5 minutes Q\&A
\item Clear visual aids and professional data visualizations
\item Effective storytelling with data and public health implications
\item Demonstration of team collaboration and individual contributions
\item Professional delivery, time management, and audience engagement
\end{itemize}

\section{Deliverables}

\subsection{Project Proposal (Week 5)}
One-page proposal including:
\begin{itemize}
\item Research question and hypothesis
\item Selected dataset and justification
\item Planned statistical approaches
\item Team member roles and responsibilities
\end{itemize}

\subsection{Final Report (Week 10)}
Comprehensive report (12-15 pages) including:
\begin{itemize}
\item Executive summary (1 page)
\item Introduction and literature review
\item Methods section with data description
\item Results with tables and figures
\item Discussion and limitations
\item Reproducible R code with comments
\item Team contribution statement
\end{itemize}

\subsection{Final Assessment Presentation Materials (Finals Week)}
\begin{itemize}
\item PowerPoint/PDF slides (20-25 slides for 20-minute presentation)
\item Live demonstration of key findings and statistical methods
\item Prepared responses for anticipated questions and statistical challenges
\item One-page executive summary handout for audience
\item Individual reflection on learning (1 page per student)
\end{itemize}

\section{Grading Rubric}

\begin{center}
\begin{tabular}{lcc}
\toprule
\textbf{Component} & \textbf{Weight} & \textbf{Points} \\
\midrule
Research Question \& Literature Review & 15\% & 15 \\
Data Management \& Exploration & 20\% & 20 \\
Statistical Analysis & 25\% & 25 \\
AI Integration \& Documentation & 15\% & 15 \\
Final Presentation \& Communication & 25\% & 25 \\
\midrule
\textbf{Total} & \textbf{100\%} & \textbf{100} \\
\bottomrule
\end{tabular}
\end{center}

\subsection{Detailed Grading Criteria}

\subsubsection{Excellent (90-100\%)}
\begin{itemize}
\item Clear, innovative research question with public health relevance
\item Comprehensive statistical analysis with appropriate methods
\item Exceptional AI integration with thorough documentation
\item Professional presentation with compelling data story
\item All team members contributed meaningfully
\end{itemize}

\subsubsection{Good (80-89\%)}
\begin{itemize}
\item Well-defined research question with clear objectives
\item Appropriate statistical methods correctly applied
\item Good AI usage with adequate documentation
\item Clear presentation with good visualizations
\item Most team members actively participated
\end{itemize}

\subsubsection{Satisfactory (70-79\%)}
\begin{itemize}
\item Basic research question addressed
\item Statistical methods mostly correct with minor errors
\item AI used but documentation incomplete
\item Adequate presentation meeting time requirements
\item Uneven team contribution evident
\end{itemize}

\subsubsection{Needs Improvement (<70\%)}
\begin{itemize}
\item Unclear research question or objectives
\item Inappropriate statistical methods or major errors
\item Minimal AI usage or poor documentation
\item Presentation lacks clarity or professionalism
\item Limited team collaboration
\end{itemize}

\section{Team Formation and Roles}

Each team of 7 students should designate:
\begin{itemize}
\item \textbf{Project Manager:} Coordinate meetings, deadlines, and submissions
\item \textbf{Data Manager:} Lead data cleaning and preprocessing
\item \textbf{Statistical Lead:} Guide analytical approach and methods
\item \textbf{AI Integration Lead:} Coordinate AI tool usage and documentation
\item \textbf{Visualization Lead:} Create figures and presentation materials
\item \textbf{Writing Lead:} Coordinate report writing and editing
\item \textbf{Presentation Lead:} Coordinate final presentation and Q\&A preparation
\end{itemize}

Note: All team members are expected to contribute to all aspects; roles are for coordination.

\section{Resources and Support}

\subsection{Technical Resources}
\begin{itemize}
\item RStudio with GitHub Copilot enabled
\item Access to OSU library databases
\item Canvas discussion board for project Q\&A
\item Sample analysis templates on Canvas
\end{itemize}

\subsection{Recommended R Packages}
\begin{verbatim}
# Data manipulation
library(tidyverse)
library(data.table)

# Statistical analysis
library(stats)
library(car)
library(lmtest)
library(tableone)

# Visualization
library(ggplot2)
library(plotly)
library(corrplot)

# Reporting
library(knitr)
library(rmarkdown)
library(stargazer)
\end{verbatim}

\section{Academic Integrity}

\begin{tcolorbox}[colback=red!5,colframe=red,boxrule=1pt,arc=2pt,left=6pt,right=6pt,top=6pt,bottom=6pt]
\textbf{Important Reminders:}
\begin{itemize}
\item All AI-generated content must be verified and attributed
\item Each team member must contribute substantively
\item Plagiarism or fabrication of data will result in project failure
\item Inter-team collaboration is encouraged, but each team's work must be unique
\item Citation of all sources (including AI tools) is required
\end{itemize}
\end{tcolorbox}

\section{Tips for Success}

\begin{enumerate}
\item \textbf{Start Early:} Data cleaning always takes longer than expected
\item \textbf{Communicate Regularly:} Schedule weekly team meetings
\item \textbf{Document Everything:} Keep detailed notes of decisions and methods
\item \textbf{Verify AI Output:} Always double-check AI-generated code and interpretations
\item \textbf{Practice Presentation:} Rehearse timing and transitions
\item \textbf{Seek Feedback:} Use office hours and peer review opportunities
\item \textbf{Focus on Story:} Ensure your analysis tells a coherent public health story
\end{enumerate}

\end{document}